\section{Readings for each weekday}
\label{sec:pc-advent-weekdays}

Start with today's date. On Sunday use \hyperref[sec:pc-advent-sundays]{the
Sunday of Advent}. On a weekday through December 16, count from the
Sunday nearest November 30 to find the current week and weekday. Check the
\hyperref[sec:pc-advent-feasts]{feasts} if it is November 30, December 8
or December 12, or if the parish calendar announces a solemnity. On a
weekday from December 17 through the morning of December 24, go straight
to \hyperref[sec:pc-advent-last-days]{the date}. After a missed day,
begin with today. If your parish reads a different authorized passage,
carry the Gospel you actually heard into prayer.

Each weekday entry gives the Mass readings in order, with the psalm or
canticle in the middle. In the Douay--Rheims Bible, 1 and 2 Samuel
appear as 1 and 2 Kings. These Advent weekday readings recur annually.
The church's weekday Year I or II makes no change to them. The
meditations, intentions and short prayers that accompany them here are
chosen for home. If time is short, read the day's Gospel and say the
prayer printed at the start of its group. If time allows, read the first
reading, give the psalm its own hearing, and spend a few minutes with
the spiritual writer beside the Gospel.

\subsection{Week one, Monday through Saturday}

The first people we meet in the weekday Gospels include a soldier asking
for his servant's life and a crowd with nothing left to eat. Keep St
Cyril of Jerusalem's \href{https://www.newadvent.org/fathers/310115.htm}{\emph{Catechetical Lecture} 15}
beside you. In §1 the child at Bethlehem and the coming King are one
Lord. In §26 the King asks what we have given him in the person of a
hungry or sick neighbor. Pray, “Shew, O Lord, thy ways to me, and teach
me thy paths” (Douay Psalm 24:4b).

\begin{description}
\item[Monday (175)] Isaiah 2:1--5. Psalm 122:1--2, 3--4b, 4cd--5,
6--7, 8--9. Matthew 8:5--11. In Year A, if Isaiah 2:1--5 was read on
Sunday, Isaiah 4:2--6 may be read instead. The centurion asks for
someone suffering under his roof. Read Cyril 15.26. Say the week's
prayer before doing one thing for a person whose care you have been given.
\item[Tuesday (176)] Isaiah 11:1--10. Psalm 72:1--2, 7--8, 12--13, 17.
Luke 10:21--24. Jesus thanks his Father for what has been disclosed to
little ones. Read Cyril 15.1 on the child wrapped for warmth. Say the
prayer with a child in mind and make room to hear what that child asks.
\item[Wednesday (177)] Isaiah 25:6--10a. Psalm 23:1--3a, 3b--4, 5, 6.
Matthew 15:29--37. Jesus knows how far the crowd must walk if he sends
them away hungry. Read Cyril 15.26 and pray before setting an extra
portion aside for someone else's supper.
\item[Thursday (178)] Isaiah 26:1--6. Psalm 118:1, 8--9, 19--21,
25--27a. Matthew 7:21, 24--27. The alternative psalm response is
Alleluia. The two builders have both heard the
words of Jesus. One of them acted on what he heard. Read Cyril 15.2
on the testing of a person's work. Pray and carry out one promise made
to a neighbor.
\item[Friday (179)] Isaiah 29:17--24. Psalm 27:1, 4, 13--14.
Matthew 9:27--31. The two blind men call after Jesus together and
follow him inside. Read Cyril 15.1. Say the prayer on behalf of someone
whose call has gone unheard in a crowd.
\item[Saturday (180)] Isaiah 30:19--21, 23--26. Psalm 147:1--2, 3--4, 5--6.
Matthew 9:35--10:1, 5a, 6--8. Jesus sees exhausted people and sends
his disciples to them. Read Cyril 15.26. Pray for the person who
will take the next shift at a bedside. See whether you can give them a break.
\end{description}

\subsection{Week two, Monday through Saturday}

St Teresa remembers keeping Christ company when she could scarcely
gather her thoughts. Spend a few minutes each day with
\href{https://ccel.org/ccel/teresa/life/life.viii.x.html}{her \emph{Life} 9.4--7}.
Pray, “Shew us, O Lord, thy mercy; and grant us thy salvation” (Douay
Psalm 84:8). When you say “us,” bring another person to Christ with you.

\begin{description}
\item[Monday (181)] Isaiah 35:1--10. Psalm 85:9ab, 10, 11--12, 13--14.
Luke 5:17--26. The men on the roof bear their friend's weight all the way
to the feet of Jesus. Read Teresa 9.4 on drawing near to someone in
affliction. Name each person now helping to carry a burden for you.
\item[Tuesday (182)] Isaiah 40:1--11. Psalm 96:1--2, 3, 10ac, 11--12, 13.
Matthew 18:12--14. Jesus lets his hearers feel the shepherd's delight
over a single sheep. Read Teresa 9.4. Pray for someone whose absence
might easily pass unnoticed. Let the person know you remember.
\item[Wednesday (183)] Isaiah 40:25--31. Psalm 103:1--2, 3--4, 8, 10.
Matthew 11:28--30. Jesus invites people who are already tired to come
as they are. Read Teresa 9.5 on staying with him as long as her thoughts
allowed. Say the week's prayer slowly enough to rest between its petitions.
\item[Thursday (184)] Isaiah 41:13--20. Psalm 145:1, 9, 10--11, 12--13ab.
Matthew 11:11--15. Isaiah promises springs for parched ground.
For Teresa, the sight of water could recall the Creator (9.6).
Say the prayer when drawing a glass of water for someone today.
\item[Friday (185)] Isaiah 48:17--19. Psalm 1:1--2, 3, 4, 6.
Matthew 11:16--19. One crowd finds fault with John's fasting and
with Jesus eating. Teresa knew how swiftly her thoughts could carry
her away (9.5--6). Where has the same habit of fault-finding robbed
you of joy in another person's good? Take that incident to prayer.
\item[Saturday (186)] Sirach 48:1--4, 9--11. Psalm 80:2ac, 3b, 15--16,
18--19. Matthew 17:9a, 10--13. The disciples come to understand that
Jesus has spoken of John, who suffered at the hands of those who failed
to recognize him. Read Teresa 9.5 on keeping company with Christ in his
suffering. Bring to prayer someone whose pain you understood only later.
\end{description}

\subsection{Week three, only through December 16}

There may be five of these weekdays or none. Use only those that arrive
before December 17. Since Thursday of week two the Gospels have been
asking who John the Baptist is and where he directs our eyes
(\emph{Introduction to the Lectionary}, no.~94). The earliest possible third Saturday is December 17,
so Saturday always takes a date. Read St John Chrysostom's
\href{https://www.newadvent.org/fathers/200136.htm}{\emph{Homily on Matthew} 36.2--3}
on the Baptist's care for his disciples. Pray, “My soul doth magnify
the Lord. And my spirit hath rejoiced in God my Saviour” (Douay Luke
1:46--47).

\begin{description}
\item[Monday (187)] Numbers 24:2--7, 15--17a. Psalm 25:4--5ab, 6, 7bc,
8--9. Matthew 21:23--27. Think of the elders discussing what an honest
answer about John might cost them. Read Chrysostom 36.2 on the disciples'
difficulty in acknowledging Jesus. Say Mary's prayer for the courage to
acknowledge good done by someone with whom you disagree.
\item[Tuesday (188)] Zephaniah 3:1--2, 9--13. Psalm 34:2--3, 6--7,
17--18, 19, 23. Matthew 21:28--32. One son changes his mind and goes
to the vineyard. Read Chrysostom 36.2. Give thanks for a chance to change
your answer, then do the work that is waiting.
\item[Wednesday (189)] Isaiah 45:6c--8, 18, 21c--25. Psalm 85:9ab, 10,
11--12, 13--14. Luke 7:18b--23. Read Chrysostom 36.2--3 on what John's
messengers see for themselves. Bring one of those people to mind while
saying Mary's prayer. Whose good news could you carry faithfully today?
\item[Thursday (190)] Isaiah 54:1--10. Psalm 30:2, 4, 5--6, 11--12a,
13b. Luke 7:24--30. Jesus waits for the messengers to leave before
speaking to the crowd about John. Read Chrysostom 36.3. Thank God by
name for a teacher whose fidelity you appreciated only after parting.
\item[Friday (191)] Isaiah 56:1--3a, 6--8. Psalm 67:2--3, 5, 7--8.
John 5:33--36. John burns like a lamp. Jesus speaks of the works the
Father has given him to do. Read Chrysostom 36.2, where John helps his
disciples look beyond himself. Thank God for the person who taught you
to make the sign of the cross or to speak the name of Jesus.
\end{description}

\subsection{December 17 to 24}
\label{sec:pc-advent-last-days}

On weekdays the date now determines the readings, even if it is still
the third week. Matthew begins with the family of Jesus and with Joseph.
Luke takes us from the angel's appearance to Zechariah to the song the
old priest sings when his son has been named. On Sunday keep the Sunday Mass. The Lectionary permits
the readings of a date displaced by Sunday to be used earlier or later
in the week, especially where a Gospel would otherwise be repeated
(rubric before no.~193). At home you may read that date's Gospel after
the Sunday Gospel or simply keep the Sunday. If December 17 is the
Third Sunday, December 24 is the Fourth. On both days follow the Sunday
entry, whatever Gospel another parish weekday might use.

St Alphonsus Liguori wrote \href{https://archive.org/details/incarnationofjes04liguuoft/page/214/mode/1up}{nine meditations for the Christmas novena}.
The first, on the Father's gift of his Son, belongs to December 16 if you
wish to begin there. His meditations II--IX accompany the dates below.
His imaginative prayer follows the love with which Christ gives himself
throughout his life. Read the numbered meditation, then answer in Mary's
words, “Behold the handmaid of the Lord: be it done to me according to
thy word” (Douay Luke 1:38).

\begin{description}
\item[December 17 (193)] Genesis 49:2, 8--10. Psalm 72:1--2, 3--4ab,
7--8, 17. Matthew 1:1--17. Read the names slowly, beginning with
Abraham. Families carry wounds as well as blessings. In meditation II
Alphonsus considers the cost to Christ of coming among us. Say Mary's
prayer for one person whose name you almost never hear at home.
\item[December 18 (194)] Jeremiah 23:5--8. Psalm 72:1--2, 12--13,
18--19. Matthew 1:18--25. The angel entrusts the child's name to
Joseph. In meditation III Alphonsus asks why God chose to come as a
child whom we might approach in love. Pray for someone to whom you have
been afraid to show tenderness.
\item[December 19 (195)] Judges 13:2--7, 24--25a. Psalm 71:3--4a,
5--6ab, 16--17. Luke 1:5--25. The angel says that Zechariah's prayer
has been heard. After long years of waiting he can scarcely believe it.
In meditation IV Alphonsus considers the obedience of the Son throughout
his life. Say Mary's prayer with one long-unanswered petition before you.
\item[December 20 (196)] Isaiah 7:10--14. Psalm 24:1--2, 3--4ab, 5--6.
Luke 1:26--38. Mary puts a question to the angel and then gives her
answer. In meditation V Alphonsus dwells on the Son's own willing
offering. Tell God what you still need to ask before saying Mary's prayer.
\item[December 21 (197)] Song of Songs 2:8--14 or Zephaniah 3:14--18a.
Psalm 33:2--3, 11--12, 20--21. Luke 1:39--45. The beloved draws near,
or Zion hears that her King is in her midst. Elizabeth recognizes her
Lord while he is still hidden. Read meditation VI, on the months Jesus
spent in Mary's womb. Remember an expectant mother far from her family.
\item[December 22 (198)] 1 Samuel 1:24--28. 1 Samuel 2:1, 4--5, 6--7,
8abcd (Hannah's canticle). Luke 1:46--56. Hear Hannah and Mary in turn.
Each remembers those whom the world overlooks. Read meditation VII on
ingratitude. Say Mary's prayer, then thank someone for a particular
kindness you have long taken for granted.
\item[December 23 (199)] Malachi 3:1--4, 23--24. Psalm 25:4--5ab,
8--9, 10, 14. Luke 1:57--66. In the Douay Bible, read Malachias
3:1--4 and 4:5--6. Elizabeth and Zechariah agree on their
son's name. The crowd wonders at the child whom God has given them.
In meditation VIII Alphonsus calls the love shown us in Christ a gift
we could never have earned. Pray for a child whose life is beginning
in difficult circumstances.
\item[December 24, morning (200)] 2 Samuel 7:1--5, 8b--12, 14a, 16.
Psalm 89:2--3, 4--5, 27, 29. Luke 1:67--79. On a weekday hear
Zechariah bless God before he speaks to his infant son. Read meditation
IX on Joseph going with Mary to Bethlehem. If you can, read the whole
canticle aloud for those who will spend tonight far from home. If today
is Sunday, use the Fourth Sunday's readings. At the evening Vigil Mass,
follow the readings for \hyperref[sec:pc-christmas]{the Nativity}.
\end{description}

\subsection{When a feast falls in Advent}
\label{sec:pc-advent-feasts}

The U.S. calendar includes the following celebrations with their own
readings. Consult the calendar of the diocese for the dedication or
patron of a particular church and other local solemnities. On an Advent
weekday before December 17, even a memorial often keeps the weekday
readings. A saint with truly proper readings has them, and the Church
also allows suitable readings for a memorial for a pastoral reason.
From December 17 to 24 the weekday normally supplies the Mass texts.
If a memorial in the General Calendar occurs, its Collect may be used
at the weekday Mass (\emph{General Instruction of the Roman Missal}
355a, 357--358). At home you may continue that day's prayer.

\begin{description}
\item[November 30. St Andrew, Apostle (684)] If November 30 falls on an
Advent weekday, the feast supplies Romans 10:9--18, Psalm 19:8, 9, 10,
11, and Matthew 4:18--22. The two responses to the same psalm draw on
Psalm 19:10 and John 6:63. If it is Sunday, keep the First Sunday of
Advent. If Advent begins in December, Andrew's feast has already passed.
Remember who first spoke the Gospel to you. Read Cyril's \emph{Lecture}
15.26 and say the first week's prayer for that person.
\item[December 8. The Immaculate Conception (689)] This is the patronal
solemnity of the United States. If December 8 is the Second Sunday of
Advent, keep that Sunday and take the solemnity on Monday, December 9.
At the solemnity read Genesis 3:9--15, 20, Psalm 98:1, 2--3ab, 3cd--4,
Ephesians 1:3--6, 11--12, and Luke 1:26--38. The Annunciation Gospel
shows us the woman whom God's grace had already prepared. With
\href{https://www.newadvent.org/fathers/0103322.htm}{St Irenaeus, \emph{Against Heresies} 3.22.4},
hear Mary's obedient answer beside Eve's disobedience. Pray Mary's words
from Luke 1:38, printed above.
\item[December 12. Our Lady of Guadalupe (690A)] On a weekday this U.S.
feast has Zechariah 2:14--17 or Revelation 11:19a and 12:1--6a, 10ab,
the canticle Judith 13:18bcde, 19, and Luke 1:26--38 or 1:39--47.
Readings from the Common of the Blessed Virgin Mary (707--712) may also
be used. In the Douay the invitation to rejoice is Zacharias 2:10--13,
and the blessing of Judith is found in 13:23--25.
When December 12 is an Advent Sunday, the U.S. calendar keeps
the Sunday. Follow the local calendar if a particular church has a
patronal solemnity. At home consider Elizabeth's welcome, with Bernard's
\emph{Homilies on the Missus est} 4.8 on Mary's answer. Pray the opening
of Mary's song (Luke 1:46--47) for visitors who will be worshipping far
from their own country.
\end{description}
